% =====================================================================
% CAPÍTULO 3: ESTADO DEL ARTE E INVESTIGACIÓN DE SOFTWARE
% Versión refactorizada con fuentes verificables
% =====================================================================

\chapter{Estado del Arte e Investigación de Software}
\label{ch:estado-arte}

% --- Ajustes editoriales locales del capítulo: sangría, espaciado y control de desbordamientos ---
\setlength{\parindent}{1.25cm}
\setlength{\parskip}{0pt}
\sloppy
\emergencystretch=2em
\raggedbottom
\section{Criterio de revisión}

La revisión del estado del arte se enfocó en herramientas reales que pudieran aportar fundamentos al diseño del aplicativo de CERMONT S.A.S. Se analizaron tres grupos: plataformas comerciales de gestión de servicios de campo o mantenimiento, repositorios abiertos relacionados con FSM/CMMS/ERP y librerías de formularios, documentos y seguridad. La finalidad de esta revisión no es copiar una solución existente, sino identificar patrones, módulos, restricciones y costos públicos que permitan justificar el desarrollo a medida.

\section{Plataformas comerciales relacionadas}

Las plataformas comerciales muestran qué funcionalidades se consideran estándar en el mercado. Sin embargo, su adopción depende de costos, capacidad de personalización, conectividad, flujo documental y ajuste al contexto de la empresa. Cuando el proveedor no publica precios o exige cotización, el libro no debe inventar valores.

\begin{table}[!htbp]
\centering
\caption{Plataformas comerciales relacionadas con gestión de campo y mantenimiento. Fuente: Elaboración propia.}
\label{tab:plataformas-comerciales}
\small
\setlength{\tabcolsep}{2pt}
\def\CommercialRow#1#2#3#4{\parbox[t]{2.55cm}{\raggedright #1} & \parbox[t]{4.1cm}{\raggedright #2} & \parbox[t]{4.55cm}{\raggedright #3} & \parbox[t]{2.75cm}{\raggedright #4}\\\midrule}
\begin{tabular}{@{}l l l l@{}}
\toprule
\textbf{Herramienta} & \textbf{Enfoque} & \textbf{Aporte al proyecto} & \textbf{Dato verificable} \\
\midrule
\CommercialRow{Jobber}{Servicios en sitio: cotización, agenda, despacho, facturación y comunicación.}{Referencia de flujo comercial para empresas de servicios.}{Planes desde USD 29 hasta USD 529 al mes \cite{jobber_pricing}.}
\CommercialRow{Odoo Field Service}{Servicio de campo en Odoo: tareas, productos, planificación y hojas de trabajo.}{Referencia para integrar servicio de campo con otros módulos.}{La documentación oficial describe tareas, productos e itinerarios \cite{odoo_field_service}.}
\CommercialRow{ServiceM8}{Contratistas: trabajos, clientes, cotizaciones y facturación.}{Referencia para movilidad y flujo técnico.}{Su página indica planes mensuales y prueba gratuita \cite{servicem8_pricing}.}
\CommercialRow{Fracttal}{Mantenimiento y gestión de activos.}{Referencia latinoamericana para CMMS y activos.}{Integra SAP, ERP, WhatsApp y hojas de cálculo \cite{fracttal}.}
\bottomrule
\end{tabular}
\end{table}

La revisión demuestra que existen soluciones maduras, pero también evidencia que una empresa contratista multiservicio puede requerir una herramienta más ajustada a sus documentos internos. El criterio de comparación no debe limitarse al precio; también debe evaluar adaptación a formatos propios, control de evidencias, operación offline, curva de aprendizaje y propiedad de los datos.

\section{Repositorios abiertos para FSM, CMMS y ERP}

\begin{table}[!htbp]
\centering
\caption{Repositorios y proyectos abiertos revisados. Fuente: Elaboración propia.}
\label{tab:repositorios-abiertos}
\small
\setlength{\tabcolsep}{2pt}
\def\RepoRow#1#2#3#4{\parbox[t]{2.45cm}{\raggedright #1} & \parbox[t]{4.05cm}{\raggedright #2} & \parbox[t]{4.0cm}{\raggedright #3} & \parbox[t]{2.55cm}{\raggedright #4}\\\midrule}
\begin{tabular}{@{}l l l l@{}}
\toprule
\textbf{Proyecto} & \textbf{Descripción} & \textbf{Relación con CERMONT S.A.S.} & \textbf{Tipo} \\
\midrule
\RepoRow{OCA Field Service}{Módulos de FSM mantenidos por OCA \cite{oca_field_service}.}{Referencia para órdenes, ubicaciones, técnicos y servicios.}{Repositorio GitHub.}
\RepoRow{ERPNext}{ERP libre con contabilidad, CRM, compras e inventario \cite{erpnext}.}{Permite comparar un ERP amplio frente a una solución focalizada.}{ERP open-source.}
\RepoRow{openMAINT}{Aplicación para activos y mantenimiento \cite{openmaint}.}{Aporta conceptos de mantenimiento e indicadores.}{CMMS/EAM.}
\RepoRow{Atlas CMMS}{CMMS autoalojable con React y Docker \cite{atlas2024}.}{Referencia para órdenes e inventario.}{Repositorio GitHub.}
\RepoRow{Liberu Maintenance}{CMMS con Laravel y Filament, gestión documental y versionamiento \cite{liberu2024}.}{Aporta patrones de documental y formularios.}{Repositorio GitHub.}
\RepoRow{FieldPro}{FSM con React 18, TypeScript y Supabase \cite{fieldpro2024}.}{Referencia por stack similar y checklist digital.}{Repositorio GitHub.}
\bottomrule
\end{tabular}
\end{table}

Estos proyectos confirman que la solución propuesta debe modelar claramente entidades como orden de trabajo, cliente, servicio, recurso, evidencia, formulario, estado y documento generado. No obstante, CERMONT S.A.S. requiere una capa adicional de adaptación a formatos operativos específicos, lo que justifica el enfoque de formularios dinámicos.

\section{Librerías para formularios dinámicos}

\begin{table}[!htbp]
\centering
\caption{Librerías de formularios dinámicos aplicables al proyecto. Fuente: Elaboración propia.}
\label{tab:librerias-formularios}
\small
\def\FormRow#1#2#3{\parbox[t]{3.4cm}{\raggedright #1} & \parbox[t]{5.0cm}{\raggedright #2} & \parbox[t]{5.7cm}{\raggedright #3}\\\midrule}
\begin{tabular}{@{}l l l@{}}
\toprule
\textbf{Librería} & \textbf{Funcionalidad relevante} & \textbf{Aplicación propuesta} \\
\midrule
\FormRow{React JSON Schema Form}{Genera formularios React de forma declarativa a partir de JSON Schema \cite{rjsf}.}{Formularios versionados si el frontend usa React.}
\FormRow{JSON Forms}{Enfoque basado en JSON Schema con soporte para React, Angular y Vue \cite{jsonforms}.}{Alternativa flexible si se requiere independencia de framework.}
\FormRow{SurveyJS}{Librería MIT para renderizar formularios JSON y enviar respuestas a una base de datos propia \cite{surveyjs}.}{Formularios largos, encuestas de inspección y checklists.}
\FormRow{Form.io}{Ecosistema de formularios JSON con componentes para frameworks web \cite{formio}.}{Evaluar si se requiere constructor visual de formularios.}
\FormRow{YAMLForms}{Generación de formularios desde esquemas YAML a PDF, HTML y DOCX \cite{yamlforms2024}.}{Referencia para línea de evolución de formularios configurables.}
\bottomrule
\end{tabular}
\end{table}

\section{Herramientas de extracción documental}

El módulo de formularios dinámicos desde documentos existentes necesita transformar entradas no estructuradas en datos editables. La Tabla \ref{tab:extraccion-documental} resume herramientas investigadas.

\begin{table}[!htbp]
\centering
\caption{Herramientas para extracción documental. Fuente: Elaboración propia.}
\label{tab:extraccion-documental}
\small
\def\ExtractRow#1#2#3{\parbox[t]{3.3cm}{\raggedright #1} & \parbox[t]{5.3cm}{\raggedright #2} & \parbox[t]{5.0cm}{\raggedright #3}\\\midrule}
\begin{tabular}{@{}l l l@{}}
\toprule
\textbf{Herramienta} & \textbf{Capacidad reportada por la fuente} & \textbf{Uso posible en el proyecto} \\
\midrule
\ExtractRow{Docling}{Procesamiento y extracción estructurada de formatos diversos, con comprensión avanzada de PDF e integración con flujos de IA \cite{docling}.}{Convertir PDF/DOCX a representación estructurada antes de generar plantillas.}
\ExtractRow{Unstructured}{Transformación de documentos complejos a formatos limpios y estructurados para modelos de lenguaje \cite{unstructured}.}{Extraer secciones, listas y bloques textuales de formatos existentes.}
\ExtractRow{PaddleOCR}{OCR y extracción estructurada para documentos PDF e imágenes, con soporte de múltiples idiomas \cite{paddleocr}.}{Procesar fotografías o PDFs escaneados de formatos de campo.}
\ExtractRow{MinerU}{Conversión de PDF, imágenes, DOCX, PPTX y XLSX en Markdown o JSON para procesamiento posterior \cite{mineru}.}{Generar entrada estructurada para un módulo futuro de plantillas dinámicas.}
\ExtractRow{PyPDFForm}{Biblioteca Python para inspeccionar y rellenar formularios PDF mediante diccionarios Python \cite{pypdfform2024}.}{Referencia para extracción y llenado de formularios PDF existentes.}
\ExtractRow{Doc2Form}{Aplicación que usa Gemini AI para convertir PDF o Word en Google Forms \cite{doc2form2024}.}{Demuestra viabilidad de extracción automatizada desde documentos no estructurados.}
\bottomrule
\end{tabular}
\end{table}

\section{Comparación de proyectos open source}

\begin{table}[!htbp]
\centering
\caption{Comparación de proyectos open source de gestión de servicios de campo. Fuente: Elaboración propia.}
\label{tab:comparacion-proyectos}
\scriptsize
\def\CompareRow#1#2#3#4#5#6{\parbox[t]{2.2cm}{\raggedright #1} & \parbox[t]{1.9cm}{\raggedright #2} & \parbox[t]{1.9cm}{\raggedright #3} & \parbox[t]{1.9cm}{\raggedright #4} & \parbox[t]{1.9cm}{\raggedright #5} & \parbox[t]{2.0cm}{\raggedright #6}\\\midrule}
\begin{tabular}{@{}l l l l l l@{}}
\toprule
\textbf{Característica} & \textbf{Atlas CMMS} & \textbf{Liberu Maint.} & \textbf{OCA Field Svc} & \textbf{FieldPro} & \textbf{CERMONT (este proyecto)} \\
\midrule
\CompareRow{Órdenes de trabajo}{Sí}{Sí}{Sí}{Sí}{Sí}
\CompareRow{Planeación de recursos}{Básica}{Sí}{Sí}{Limitada}{Sí (kits típicos)}
\CompareRow{Checklists digitales}{No}{No}{No}{Sí (48 ítems)}{Sí}
\CompareRow{Evidencias fotográficas}{No}{No}{Limitado}{Sí}{Sí}
\CompareRow{Generación de PDF}{Sí}{Sí}{Sí}{Sí (facturas, reportes)}{Sí (informes, actas)}
\CompareRow{Operación offline}{No}{No}{Limitado}{Sí (PWA)}{Sí (PWA + IndexedDB)}
\CompareRow{Gestión documental}{Básica}{Sí (versionado)}{Limitada}{No}{Sí (trazabilidad completa)}
\CompareRow{Formularios personalizables}{No}{Sí}{No}{No}{Propuesto (futuro)}
\CompareRow{Cierre administrativo}{No}{No}{No}{No}{Sí (SES, facturación)}
\CompareRow{Stack principal}{React + Docker}{Laravel + PHP}{Python (Odoo)}{React + Supabase}{Next.js + Express + MongoDB}
\CompareRow{Licencia}{AGPLv3 / Com.}{MIT}{AGPLv3}{Propietaria}{Académica / Propietaria}
\bottomrule
\end{tabular}

\vspace{0.3cm}
\raggedright
\footnotesize
\textbf{Fuente:} Elaboración del autor con base en la revisión de repositorios GitHub, documentación oficial de cada proyecto y análisis comparativo de funcionalidades. Las celdas marcadas como ``No'' indican ausencia de la funcionalidad como módulo nativo; algunas pueden estar disponibles mediante extensiones o configuraciones adicionales no cubiertas en esta revisión.
\end{table}

\section{Vacíos identificados y posicionamiento del proyecto}

La investigación muestra que las plataformas comerciales resuelven partes del problema, pero no necesariamente absorben los documentos existentes de una empresa contratista ni reflejan sus formatos internos sin parametrización o desarrollo adicional. Los repositorios abiertos ofrecen bases técnicas, pero requieren adaptación funcional y soporte de implementación. Las librerías de formularios y extracción documental permiten plantear una arquitectura más flexible que los formularios codificados manualmente.

El proyecto se posiciona como una solución aplicada, enfocada en el flujo real de CERMONT S.A.S.: planeación, ejecución, evidencias, informes, actas y cierre administrativo. Su aporte no está en prometer ahorros no medidos, sino en construir un artefacto de software con base documental, justificar técnicamente sus módulos y definir una validación que pueda demostrar resultados cuando existan datos reales.

% =====================================================================
\section{Análisis de madurez del mercado FSM}
\label{sec:madurez-fsm-mercado}

El mercado de software de gestión de servicios de campo (FSM) ha experimentado una evolución significativa en la última década. Según el informe Magic Quadrant for Field Service Management de Gartner \cite{gartner2024}, el mercado se consolida alrededor de capacidades como la inteligencia artificial para programación predictiva, la realidad aumentada para asistencia remota y la analítica avanzada para optimización de fuerza laboral. Los líderes del cuadrante (ServiceMax, IFS, Oracle, SAP) ofrecen plataformas empresariales con amplia cobertura funcional, pero con costos de licencia que pueden oscilar entre USD 10,000 y USD 50,000 anuales para implementaciones en empresas medianas, según el mismo informe.

Un estudio de Aberdeen Group \cite{aberdeen2019} encontró que las empresas con mejores prácticas en FSM logran una tasa de resolución en primera visita (First-Time Fix Rate) del 77\%, comparada con el 63\% del promedio de la industria, atribuyendo esta diferencia a la disponibilidad de información completa en el punto de servicio. Este hallazgo refuerza la importancia de que el aplicativo de CERMONT S.A.S. proporcione a los técnicos de campo acceso a toda la información relevante (historial del activo, checklists, kits de recursos) antes y durante la ejecución.

En el contexto latinoamericano, la adopción de herramientas FSM presenta patrones diferenciados. Las empresas multinacionales que operan en la región tienden a utilizar las plataformas corporativas globales (SAP, Oracle, IFS), mientras que las PYMEs contratistas, como CERMONT S.A.S., enfrentan barreras de adopción relacionadas con el costo de licenciamiento, la complejidad de implementación y la capacidad limitada de personalización de las soluciones comerciales. Este segmento de mercado representa una oportunidad para soluciones desarrolladas a medida que, aunque no compiten en amplitud funcional con las plataformas empresariales, ofrecen un ajuste preciso a los procesos documentales existentes de la organización.

\section{Análisis detallado de repositorios open source}
\label{sec:deep-dive-repos}

A continuación se presenta un análisis en profundidad de tres repositorios que resultaron particularmente relevantes para el diseño del aplicativo de CERMONT S.A.S., ya sea por su stack tecnológico, su arquitectura de módulos o sus funcionalidades afines.

\subsection{FieldPro: referencia de stack y funcionalidades PWA}

FieldPro \cite{fieldpro2024} es una plataforma de gestión de servicios de campo construida con React 18, TypeScript, Vite 6, Tailwind CSS y Supabase (PostgreSQL + Row Level Security). Su arquitectura resultó particularmente instructiva para el proyecto CERMONT por tres razones. Primera, implementa un checklist digital de 48 ítems con estado tri-estado (completado, pendiente, no aplica), que sirvió como referencia para el diseño de los checklists de inspección del aplicativo. Segunda, incorpora soporte PWA con caché offline, demostrando la viabilidad técnica de la operación desconectada en aplicaciones web modernas. Tercera, su enfoque de generación de documentos PDF (facturas, reportes de servicio, cotizaciones) a partir de datos estructurados validó la decisión de utilizar \texttt{pdf-lib} como motor de generación documental en el proyecto.

Sin embargo, FieldPro presenta limitaciones que justifican el desarrollo de una solución propia para CERMONT S.A.S. Carece de un módulo de cierre administrativo que conecte la ejecución técnica con la facturación y el pago, no soporta formularios dinámicos configurables por tipo de servicio, y su modelo de datos está orientado a empresas de servicios de equipos más que a contratistas multiservicio con flujos documentales complejos.

\subsection{OCA Field Service: patrón de descomposición modular}

El conjunto de módulos Field Service de la Odoo Community Association \cite{oca_field_service}, con 179 estrellas y 90 contribuidores en GitHub, representa uno de los proyectos open source más maduros en el dominio FSM. Su arquitectura ejemplifica un patrón de descomposición modular que fue adoptado en el diseño del aplicativo CERMONT: cada funcionalidad se organiza como un módulo independiente (territorios, órdenes, CRM, portal, proyectos, inventario, vehículos) con dependencias explícitas entre ellos. Esta estructura evita el acoplamiento excesivo y permite que cada módulo evolucione a su propio ritmo.

La principal diferencia con el proyecto CERMONT radica en la dependencia del ecosistema Odoo. OCA Field Service hereda la infraestructura de Odoo (ORM, sistema de permisos, motor de reportes, framework web), lo que limita su portabilidad a entornos que no utilicen Odoo como ERP base. El aplicativo CERMONT, en contraste, se diseñó como una plataforma autónoma que puede operar sin depender de un ERP específico, comunicándose mediante APIs REST estándar.

\subsection{Atlas CMMS: referencia de despliegue autoalojado}

Atlas CMMS \cite{atlas2024} es un sistema de gestión de mantenimiento computarizado autoalojado, desarrollado con React para el frontend y desplegable mediante Docker Compose. Su arquitectura de despliegue basada en contenedores demostró la viabilidad de ofrecer una solución CMMS que el cliente puede instalar en su propia infraestructura, un modelo relevante para CERMONT S.A.S. considerando que sus datos operativos y financieros deben permanecer bajo su control.

Las funcionalidades de Atlas CMMS —creación y asignación de órdenes de trabajo, registro de tiempos, automatización de órdenes mediante disparadores, exportación de reportes— cubren adecuadamente el dominio de mantenimiento de activos fijos en instalaciones industriales. Sin embargo, su alcance no abarca la gestión integral de servicios contratados para terceros, donde la trazabilidad documental, la generación de actas e informes para clientes y el cierre administrativo con facturación constituyen requisitos críticos que diferencian el proyecto CERMONT de un CMMS convencional.

\section{Antecedentes investigativos en el contexto colombiano}
\label{sec:antecedentes-colombia}

La revisión de trabajos de grado en repositorios institucionales de universidades colombianas permitió identificar tres investigaciones relacionadas con la digitalización de procesos técnicos mediante aplicaciones web. Estos antecedentes, aunque abordan dominios diferentes al de CERMONT S.A.S., proporcionan evidencia del interés académico en el tema y permiten situar el presente trabajo dentro de una línea de investigación activa en la ingeniería colombiana.

\begin{table}[!htbp]
\centering
\caption{Antecedentes académicos colombianos relacionados con la problemática. Fuente: Elaboración propia.}
\label{tab:antecedentes-colombia}
\small
\setlength{\tabcolsep}{4pt}
\def\AntecedentRow#1#2#3#4{\parbox[t]{2.4cm}{\raggedright #1} & \parbox[t]{3.4cm}{\raggedright #2} & \parbox[t]{3.8cm}{\raggedright #3} & \parbox[t]{4.0cm}{\raggedright #4}\\\midrule}
\begin{tabular}{@{}l l l l@{}}
\toprule
\textbf{Autor y año} & \textbf{Enfoque general} & \textbf{Aporte relevante} & \textbf{Diferencia con este proyecto} \\
\midrule
\AntecedentRow{García López (2021) \cite{garcia2021}}{Sistema web para mantenimiento de flota urbana.}{Demuestra viabilidad de digitalizar controles operativos con tecnologías web.}{Enfocado en activos propios; no cubre cierre administrativo con facturación.}
\AntecedentRow{Rodríguez y Torres (2022) \cite{rodriguez2022}}{Aplicación móvil para órdenes de trabajo en telecomunicaciones.}{Muestra interés por movilidad y registro en campo.}{Prioriza movilidad; no incluye trazabilidad de cierre documental.}
\AntecedentRow{Martínez Suárez (2020) \cite{martinez2020}}{Sistema web para mantenimiento preventivo en PYME industrial.}{Demuestra control de actividades técnicas con base documental.}{Alcance limitado a mantenimiento; no cubre servicios a terceros ni facturación.}
\bottomrule
\end{tabular}
\end{table}

Estos antecedentes confirman que la digitalización de procesos de mantenimiento y servicios técnicos mediante aplicaciones web es un área de interés activo en la ingeniería colombiana. Sin embargo, la revisión también evidencia que los trabajos existentes se han concentrado en mantenimiento de activos propios (flotas, equipos industriales, infraestructura de telecomunicaciones) más que en la gestión integral de servicios contratados para terceros, donde la trazabilidad documental, la generación de actas e informes y el cierre administrativo constituyen requisitos críticos que diferencian el presente proyecto.

\section{Tendencias tecnológicas emergentes y su relación con el proyecto}
\label{sec:tendencias-emergentes}

El desarrollo del aplicativo CERMONT se sitúa en la intersección de varias tendencias tecnológicas que están transformando la gestión de servicios de campo. A continuación se analizan tres tendencias particularmente relevantes y su relación con las decisiones de diseño adoptadas.

\textbf{Operación offline-first.} La combinación de Progressive Web Apps (PWA) con almacenamiento local IndexedDB y sincronización diferida se ha consolidado como el enfoque predominante para aplicaciones de campo en zonas con conectividad intermitente \cite{gartner2024}. El proyecto CERMONT adopta este enfoque mediante Service Workers que interceptan peticiones HTTP y almacenan datos localmente, con un motor de sincronización que aplica retroceso exponencial y deduplicación mediante \texttt{clientMutationId}. Esta arquitectura permite que los técnicos de campo capturen evidencias y completen checklists independientemente de la disponibilidad de red, una capacidad crítica para operaciones en zonas rurales de Arauca.

\textbf{Formularios dinámicos configurables.} La tendencia hacia formularios definidos mediante esquemas declarativos (JSON Schema, YAML) en lugar de componentes de interfaz codificados manualmente responde a la necesidad de adaptabilidad en entornos multiservicio \cite{rjsf,jsonforms}. El proyecto CERMONT incorpora esta tendencia mediante la especificación de formularios como plantillas versionadas en \texttt{@cermont/shared-types}, con un motor de renderizado que genera la interfaz de captura a partir del esquema. La evolución futura hacia formularios extraídos automáticamente de documentos PDF/WORD existentes representa la convergencia de esta tendencia con las capacidades emergentes de extracción documental.

\textbf{Generación documental automatizada.} La capacidad de generar documentos PDF, DOCX y otros formatos a partir de datos estructurados almacenados en el sistema elimina la necesidad de recaptura manual y garantiza consistencia entre los datos operativos y los documentos entregados al cliente \cite{pypdfform2024}. El proyecto CERMONT implementa esta capacidad mediante \texttt{pdf-lib} para la generación de informes técnicos y actas de entrega, con plantillas que insertan automáticamente datos de la orden, resultados de checklists y evidencias fotográficas.

\section{Justificación de la solución a medida frente a alternativas existentes}
\label{sec:justificacion-a-medida}

La revisión del estado del arte permite articular una justificación fundamentada para el desarrollo de una solución a medida en lugar de adoptar una plataforma existente. Esta justificación se estructura en cinco criterios de evaluación, presentados en la Tabla \ref{tab:criterios-evaluacion}.

\begin{table}[!htbp]
\centering
\caption{Criterios de evaluación para la decisión de desarrollar vs. adoptar. Fuente: Elaboración propia.}
\label{tab:criterios-evaluacion}
\scriptsize
\setlength{\tabcolsep}{2pt}
\def\CriteriaRow#1#2#3#4#5{\parbox[t]{2.55cm}{\raggedright #1} & \parbox[t]{3.7cm}{\raggedright #2} & \parbox[t]{2.05cm}{\raggedright #3} & \parbox[t]{2.05cm}{\raggedright #4} & \parbox[t]{2.05cm}{\raggedright #5}\\\midrule}
\begin{tabular}{@{}l l l l l@{}}
\toprule
\textbf{Criterio} & \textbf{Descripción} & \textbf{Plataformas comerciales} & \textbf{Open source} & \textbf{Desarrollo a medida} \\
\midrule
\CriteriaRow{Adaptación a formatos propios}{Ajuste a formatos reales de CERMONT.}{Baja}{Media}{Alta}
\CriteriaRow{Costo total de propiedad}{Costo acumulado a tres años.}{Alto}{Medio}{Bajo}
\CriteriaRow{Independencia tecnológica}{Control del código y de los datos.}{Baja}{Alta}{Alta}
\CriteriaRow{Operación offline}{Funcionamiento sin conectividad.}{Variable}{Variable}{Alta}
\CriteriaRow{Evolución funcional}{Incorporar y modificar módulos.}{Limitada}{Media}{Alta}
\bottomrule
\end{tabular}
\end{table}

El análisis de estos criterios confirma que, para el contexto específico de CERMONT S.A.S. —una empresa contratista multiservicio con formatos operativos propios, operación en zonas de conectividad variable y un volumen de órdenes que no justifica la inversión en plataformas empresariales—, el desarrollo a medida constituye la alternativa técnicamente más adecuada. Esta conclusión no desconoce el valor de las plataformas existentes, que ofrecen madurez funcional y soporte profesional; simplemente reconoce que el costo de adaptar dichas plataformas a la realidad documental y operativa de CERMONT S.A.S. superaría el costo de construir una solución focalizada.

% =====================================================================
\section{Análisis de costos del ecosistema de software}
\label{sec:costos-ecosistema}

Un factor determinante en la decisión de desarrollar una solución a medida es el costo de las alternativas existentes. A continuación se presenta un análisis de los costos públicamente disponibles de las plataformas comerciales revisadas, con el propósito de proporcionar un contexto económico verificable a la decisión técnica. Los valores presentados provienen exclusivamente de páginas oficiales de los proveedores, consultadas durante la investigación.

\textbf{Jobber} publica planes de suscripción escalonados según funcionalidad \cite{jobber_pricing}. El plan básico (\textit{Lite}), orientado a negocios unipersonales, tiene un costo de USD 29 por mes facturado anualmente. El plan intermedio (\textit{Connect}), que incluye programación avanzada y recordatorios automáticos, cuesta USD 119 por mes. El plan avanzado (\textit{Grow}), con informes financieros y seguimiento de gastos, alcanza USD 249 por mes. El plan empresarial (\textit{Scale}), que incorpora múltiples ubicaciones y reportes personalizados, llega a USD 529 por mes. Para una empresa con el perfil operativo de CERMONT S.A.S., el plan relevante sería al menos \textit{Grow}, lo que representa un costo anual de aproximadamente USD 2,988, sin incluir costos de implementación, migración de datos históricos ni capacitación del personal.

\textbf{Odoo} ofrece un modelo de precios diferenciado según el número de usuarios y las aplicaciones contratadas \cite{odoo_pricing}. El módulo de \textit{Field Service} forma parte del ecosistema de aplicaciones de Odoo, cuyo costo varía según la cantidad de usuarios y la inclusión de otros módulos empresariales (contabilidad, inventario, CRM). La documentación oficial de Field Service \cite{odoo_field_service} describe funcionalidades que, si bien cubren parcialmente las necesidades de gestión de servicios de campo, requerirían personalización adicional para adaptarse a los formatos operativos propios de CERMONT S.A.S.

\textbf{ServiceM8} ofrece planes desde aproximadamente USD 29 por mes para funcionalidades básicas, con progresión de precios según el número de usuarios y las capacidades requeridas \cite{servicem8_pricing}. Su enfoque en contratistas y negocios de servicios lo hace particularmente relevante como referencia, aunque su modelo de datos está optimizado para servicios de corta duración más que para proyectos complejos con flujos documentales extensos.

\textbf{Fracttal}, con presencia en el mercado latinoamericano, no publica precios estándar en su sitio web, requiriendo contacto directo para cotización. Esta práctica, común en plataformas empresariales, dificulta la comparación objetiva de costos y refuerza el argumento a favor de una solución con costos transparentes y controlables \cite{fracttal}.

El análisis de costos, sumado a los criterios cualitativos de la Tabla \ref{tab:criterios-evaluacion}, confirma que la inversión requerida para adoptar una plataforma comercial —incluyendo licencias, implementación, migración de datos, capacitación y personalización— no se justifica para el volumen actual de operaciones de CERMONT S.A.S., particularmente cuando dicha inversión no garantiza la adaptación a los formatos operativos propios de la empresa. El desarrollo a medida, cuyo costo principal es el tiempo de ingeniería del practicante, produce un sistema ajustado a los procesos reales de la organización, con independencia tecnológica total y sin costos recurrentes de licenciamiento.

% --- AMPLIACIÓN ACADÉMICA INTEGRADA CAPÍTULO 3 ---

\section{Criterios académicos para comparar soluciones existentes}
\label{sec:criterios-comparacion-soluciones}

La comparación con herramientas comerciales y repositorios abiertos no debe limitarse a enumerar nombres de software. Para que el análisis tenga valor académico, cada alternativa debe examinarse según criterios vinculados al problema de CERMONT S.A.S.: cobertura del flujo desde solicitud hasta pago, capacidad de configurar formularios, soporte para evidencias, posibilidad de operación móvil/offline, costos de adopción, facilidad de personalización, integración documental y control de cierre administrativo.

\begin{table}[!htbp]
\centering
\scriptsize
\caption{Criterios de comparación utilizados para evaluar alternativas frente al aplicativo a medida. Fuente: elaboración propia.}
\label{tab:criterios-comparacion-alternativas}
\begin{adjustbox}{max width=\textwidth}
\begin{tabular}{p{3.2cm}p{5.7cm}p{5.7cm}}
\toprule
\textbf{Criterio} & \textbf{Pregunta de evaluación} & \textbf{Relevancia para CERMONT S.A.S.}\\
\midrule
Cobertura del flujo & ¿La herramienta cubre solicitud, propuesta, PO, planeación, ejecución, informes, actas, SES, factura y pago? & El proceso real de CERMONT no termina en la ejecución técnica, sino en el cierre administrativo.\\
Configurabilidad documental & ¿Permite adaptar formularios y plantillas a formatos propios de la empresa? & La empresa usa documentos diversos para CCTV, líneas de vida, planeación y soporte HES.\\
Operación móvil y offline & ¿Permite trabajar con conectividad intermitente sin perder datos? & El registro de evidencias y checklists se realiza en campo.\\
Trazabilidad y auditoría & ¿Conserva historial, estados, responsables y evidencias por orden? & La trazabilidad permite reconstruir el avance de cada servicio.\\
Costo y adopción & ¿El costo y la complejidad de implementación son viables para la empresa? & Las soluciones robustas pueden ser costosas o rígidas para una contratista multiservicio local.\\
Control de costos reales & ¿Relaciona presupuesto, ejecución y facturación? & La empresa necesita comparar costos reales frente a lo estimado.\\
\bottomrule
\end{tabular}
\end{adjustbox}
\end{table}

\section{Posicionamiento del aplicativo frente a FSM, CMMS y ERP}
\label{sec:posicionamiento-fsm-cmms-erp}

Las herramientas FSM se especializan en la gestión de técnicos, rutas, órdenes y atención en campo. Los sistemas CMMS se orientan a mantenimiento de activos, programación de intervenciones e historial de equipos. Los ERP integran procesos empresariales amplios, pero suelen requerir parametrización compleja, costos de implementación y adaptación organizacional significativa. El aplicativo propuesto para CERMONT S.A.S. se ubica en un punto intermedio: toma elementos de FSM para la ejecución en campo, elementos de CMMS para control de actividades técnicas y elementos de ERP para seguimiento administrativo, pero los adapta a un flujo propio de 14 pasos.

Este posicionamiento justifica el desarrollo a medida. CERMONT requiere una herramienta que no solo asigne órdenes a técnicos, sino que conecte propuesta, PO, planeación, evidencias, informe, acta, SES, factura, pago y costos reales. La mayoría de soluciones generalistas obliga a modificar el proceso de la empresa para encajar en la herramienta; el proyecto busca lo contrario: modelar la herramienta desde el proceso real y, gradualmente, introducir estandarización mediante reglas, plantillas y estados.

\section{Vacíos del mercado identificados para el caso CERMONT}
\label{sec:vacios-mercado-cermont}

El estado del arte permite identificar cuatro vacíos principales. El primero es la dificultad de adaptar herramientas comerciales a formatos internos sin incurrir en parametrizaciones extensas. El segundo es la separación frecuente entre operación técnica y cierre administrativo: muchas herramientas gestionan órdenes y evidencias, pero no necesariamente modelan SES, facturación, aprobación y pago dentro del mismo pipeline. El tercero es la falta de un enfoque explícito sobre costos reales comparados con propuesta inicial para cada servicio. El cuarto es la necesidad de operación offline y de sincronización controlada para escenarios de campo.

El aplicativo propuesto se justifica porque esos vacíos se reflejan directamente en las fallas documentadas por CERMONT. No basta con adquirir una herramienta que administre órdenes; se requiere una solución que traduzca la lógica del proceso interno en módulos, estados y validaciones. Por esta razón, el desarrollo a medida se presenta como una alternativa académicamente defendible, siempre que el libro distinga con claridad entre lo implementado, lo parcialmente implementado y lo propuesto como evolución futura.

\section{Criterio de uso de precios y costos externos}
\label{sec:criterio-precios-costos-externos}

Los costos de herramientas comerciales solo deben usarse cuando provengan de fuentes públicas verificables y con fecha de consulta. En ningún caso deben presentarse como costos propios de CERMONT S.A.S. si la empresa no ha entregado documentos financieros que lo respalden. Por tanto, la comparación económica debe limitarse a costos de mercado, licenciamiento o suscripción publicados por proveedores, y su función debe ser contextualizar la decisión de desarrollar a medida, no demostrar ahorros internos no medidos.

En el libro, cualquier mención a ahorro, retorno de inversión, reducción porcentual de tiempos o mejora cuantitativa debe quedar condicionada a evidencia de validación. Si no existen mediciones reales, se debe escribir que el indicador se propone para una fase piloto posterior. Esta regla conserva la integridad académica del documento y evita afirmaciones que puedan ser cuestionadas durante la sustentación.
